\usepackage{iftex}
\ifLuaTeX\else
  \PackageError{rlspine}{This book requires LuaLaTeX}{Run lualatex book.tex.}
\fi

\usepackage[paperwidth=7in,paperheight=10in,inner=0.78in,outer=0.82in,top=0.76in,bottom=0.82in,headsep=0.18in,footskip=0.34in]{geometry}
\usepackage{xcolor}
\definecolor{ink}{HTML}{17202A}
\definecolor{navy}{HTML}{183B56}
\definecolor{teal}{HTML}{1F6F78}
\definecolor{gold}{HTML}{C8872D}
\definecolor{mist}{HTML}{EEF4F5}
\definecolor{sand}{HTML}{F7F2E8}
\definecolor{rose}{HTML}{F7ECE9}
\color{ink}

\usepackage{fontspec}
% Prefer EB Garamond when installed; fall back to TeX Gyre Pagella, which every TeX Live ships.
\IfFontExistsTF{EBGaramond-Regular.otf}{%
  \setmainfont{EBGaramond-Regular.otf}[BoldFont=EBGaramond-Bold.otf,ItalicFont=EBGaramond-Italic.otf,BoldItalicFont=EBGaramond-BoldItalic.otf,Ligatures={TeX,Common},Numbers={OldStyle,Proportional}]%
  \newfontfamily\liningfont{EBGaramond-Regular.otf}[BoldFont=EBGaramond-Bold.otf,ItalicFont=EBGaramond-Italic.otf,Numbers={Lining,Monospaced}]%
}{%
  \setmainfont{texgyrepagella-regular.otf}[BoldFont=texgyrepagella-bold.otf,ItalicFont=texgyrepagella-italic.otf,BoldItalicFont=texgyrepagella-bolditalic.otf,Ligatures=TeX]%
  \newfontfamily\liningfont{texgyrepagella-regular.otf}[BoldFont=texgyrepagella-bold.otf,ItalicFont=texgyrepagella-italic.otf]%
}
\usepackage[tracking=true,protrusion=true,expansion=false]{microtype}

\usepackage{amsmath,amssymb,amsthm,mathtools,bm}
\DeclareMathOperator{\E}{\mathbb{E}}
\DeclareMathOperator{\Var}{Var}
\DeclareMathOperator{\Cov}{Cov}
\DeclareMathOperator*{\argmax}{arg\,max}
\DeclareMathOperator*{\argmin}{arg\,min}
\newcommand{\R}{\mathbb{R}}
\newcommand{\Prob}{\mathbb{P}}
\newcommand{\given}{\mathrel{|}}
\newcommand{\1}{\mathbf{1}}
\newcommand{\norm}[1]{\left\lVert #1\right\rVert}
\newcommand{\abs}[1]{\left\lvert #1\right\rvert}
\newcommand{\inner}[2]{\left\langle #1,#2\right\rangle}

\usepackage{booktabs,tabularx,array}
\usepackage{comment}
% The EPUB build reads these blocks; the print build deliberately discards them.
\excludecomment{epubonly}
\renewcommand{\arraystretch}{1.12}
% Ragged-right X columns; pandoc's LaTeX reader understands X but not custom column types.
\renewcommand{\tabularxcolumn}[1]{>{\raggedright\arraybackslash}p{#1}}
\usepackage{enumitem}
\setlist[itemize]{leftmargin=1.25em,itemsep=0.18em,topsep=0.35em}
\setlist[enumerate]{leftmargin=1.55em,itemsep=0.2em,topsep=0.35em}

\usepackage[most]{tcolorbox}
\usetikzlibrary{arrows.meta,positioning,calc}
\tcbset{enhanced,boxrule=0.45pt,arc=1.2mm,left=2.2mm,right=2.2mm,top=1.7mm,bottom=1.7mm,before skip=7pt,after skip=7pt}
\newtcolorbox{mentalmodel}{colback=mist,colframe=teal,title={Mental model},fonttitle=\bfseries}
\newtcolorbox{failuremode}{colback=rose,colframe=gold,title={Failure mode},fonttitle=\bfseries}
\newtcolorbox{checkpoint}{colback=sand,colframe=navy,title={Checkpoint},fonttitle=\bfseries}
\newtcolorbox{takeaway}{colback=navy!5,colframe=navy,title={What survives},fonttitle=\bfseries}
\newtcolorbox{algorithm}[1]{colback=white,colframe=teal,title={#1},fonttitle=\bfseries}

\theoremstyle{plain}
\newtheorem{theorem}{Theorem}[chapter]
\newtheorem{proposition}[theorem]{Proposition}
\newtheorem{lemma}[theorem]{Lemma}
\theoremstyle{definition}
\newtheorem{definition}[theorem]{Definition}
\newtheorem{example}[theorem]{Example}

\usepackage{fancyhdr}
\pagestyle{fancy}
\fancyhf{}
\fancyhead[LE]{\small\itshape\leftmark}
\fancyhead[RO]{\small\itshape\rightmark}
\fancyfoot[C]{\small\liningfont\thepage}
\renewcommand{\headrulewidth}{0.25pt}
\renewcommand{\footrulewidth}{0pt}
\renewcommand{\chaptermark}[1]{\markboth{#1}{}}
\renewcommand{\sectionmark}[1]{\markright{\thesection\quad #1}}
\fancypagestyle{plain}{\fancyhf{}\fancyfoot[C]{\small\liningfont\thepage}\renewcommand{\headrulewidth}{0pt}}

\makeatletter
\def\@makechapterhead#1{%
  \vspace*{20pt}%
  {\parindent\z@\raggedright\normalfont
   \ifnum\c@secnumdepth>\m@ne
     {\small\color{teal}\textsc{Chapter \thechapter}}\par\vskip 8pt
   \fi
   {\Huge\color{navy}#1\par}\vskip 9pt
   {\color{gold}\rule{0.23\linewidth}{1.2pt}}\par\vskip 23pt}}
\def\@makeschapterhead#1{%
  \vspace*{20pt}%
  {\parindent\z@\raggedright\normalfont
   {\Huge\color{navy}#1\par}\vskip 9pt
   {\color{gold}\rule{0.23\linewidth}{1.2pt}}\par\vskip 23pt}}
\makeatother

\usepackage[hidelinks,unicode]{hyperref}
\hypersetup{
  colorlinks=true,
  linkcolor=navy,
  citecolor=teal,
  urlcolor=teal,
  pdftitle={BOOK TITLE},
  pdfauthor={Prepared for LEARNER},
  pdfsubject={One-line subtitle},
  pdfkeywords={topic one, topic two}
}

\setlength{\parindent}{1.15em}
\setlength{\parskip}{0pt}
\linespread{1.04}
\raggedbottom
\widowpenalty=10000
\clubpenalty=10000
\displaywidowpenalty=10000
\emergencystretch=2em
\allowdisplaybreaks[2]

\newcommand{\term}[1]{\textbf{#1}}
\newcommand{\why}{\paragraph{Why this is true.}}
\newcommand{\practice}[1]{\begin{checkpoint}#1\end{checkpoint}}
\newenvironment{readingnote}
  {\par\smallskip\noindent\begingroup\small\color{teal}}
  {\par\endgroup\smallskip}
\newenvironment{intuitionnote}
  {\par\smallskip\noindent\begingroup\color{navy}}
  {\par\endgroup\smallskip}
\newcommand{\reading}[1]{\begin{readingnote}\textsc{Read deeper:} #1\end{readingnote}}
\newcommand{\intuition}[1]{\begin{intuitionnote}\textbf{Picture it.}~#1\end{intuitionnote}}
